\section{Conclusion}
\label{sec:conclusion}

Spectrum-Matched Circuit Design turns the choice of a quantum circuit for a PDE into
arithmetic: the operator's Fourier symbol and the forcing give the frequencies the
solution needs, and the encoding scalings, depth, qubit count and entangler follow by
formula. The circuit it builds produces exactly the designed frequencies, and those
frequencies matter: against the same circuit with random ones, the designed circuit is
$6.5\times$ more accurate on the heat equation, $236\times$ on Poisson and $9.9\times$ on
groundwater, on every seed.

The explainability instruments showed what a single circuit cannot do. All of its
Fourier amplitudes depend on the same few angles, so on a solution that mixes a slow and
a fast wave the optimiser trades one for the other. Mixing eight copies of the designed
circuit through a linear head keeps the frequency set and frees the amplitudes. On
waterlogging of farmland from leaking canals, that model reaches 0.09\% error, $18\times$
more accurate than a classical network of the same size, and answers the planning
question to within 1\,m.

The results also mark the method's limits. On a smoothing operator no model beats a
plain MLP, as SMCD predicts before training; on Poisson a classical model handed the
same frequencies is $3.3\times$ more accurate than the circuit; and at high Helmholtz
wavenumbers every model fails. The most direct next steps are running the designed
circuits on quantum hardware, extending SMCD to problems of higher dimension, and
understanding when a trainable encoder lets the circuit beat fixed classical features at
the same frequencies, as it does on the groundwater problem.
